\section{紧致空间}\label{section:ch03-compact-spaces}

紧致性远不及连通性那样自然。在拓扑学的最初阶段，人们就已经注意到实直线上闭区间 $[a,b]$ 具有一种特性，它对于证明极大值定理和一致连续性定理等结论起着至关重要的作用。但对于在任意拓扑空间中如何表述这个特性，人们长期不得而知。起初，人们以为 $[a,b]$ 的这一特性所指的是 $[a,b]$ 中任何一个无穷子集都有极限点，并且将其尊称为紧致性。后来，数学家们才意识到这种提法并未触及问题的本质，而藉助空间的开覆盖所给出的一个较强的提法更为恰切。后面这种提法就是我们现在所讲的紧致性。它不像前者那样自然或直观，在展示其效用之前我们需要先来熟悉它。

\begin{definition}\label{definition:ch03-covering}
设 $\mathcal{A}$ 是空间 $X$ 的一个子集族，如果 $\mathcal{A}$ 的成员之并等于 $X$，则称 $\mathcal{A}$ \textbf{覆盖} $X$，或者称 $\mathcal{A}$ 是 $X$ 的一个\textbf{覆盖}（covering）。如果 $\mathcal{A}$ 的每一成员都是 $X$ 的开子集，则称它为 $X$ 的一个\textbf{开覆盖}（open covering）。
\end{definition}

\begin{definition}\label{definition:ch03-compact}
若 $X$ 的任何一个开覆盖 $\mathcal{A}$，包含着一个覆盖 $X$ 的有限子族，则称空间 $X$ 是\textbf{紧致}（compact）的。
\end{definition}

\begin{example}\label{example:ch03-real-not-compact}
实直线 $\mathbb{R}$ 不是紧致的，因为由开区间
\begin{align*}
\mathcal{A}=\{\,(n,n+2)\mid n\in\mathbb{Z}\,\}
\end{align*}
所组成的 $\mathbb{R}$ 的覆盖并不包含覆盖 $\mathbb{R}$ 的任何有限子族。
\end{example}

\begin{example}\label{example:ch03-convergent-seq-compact}
$\mathbb{R}$ 的子空间
\begin{align*}
X=\{0\}\cup\{\,1/n\mid n\in\mathbb{Z}_{+}\,\}
\end{align*}
是紧致的。任意给定 $X$ 的一个开覆盖 $\mathcal{A}$，$\mathcal{A}$ 中总有一个成员 $U$ 包含 $0$，那么集合 $U$ 除了有限多个点 $1/n$ 外，包含着 $X$ 的所有其余点。对于 $X$ 中每一个 $U$ 以外的点，选取 $\mathcal{A}$ 中包含它的一个成员。于是 $\mathcal{A}$ 中这些成员连同成员 $U$ 便组成了 $\mathcal{A}$ 的一个覆盖 $X$ 的有限子族。
\end{example}

\begin{example}\label{example:ch03-finite-space-compact}
任何一个仅含有有限多个点的空间必然是紧致的，因为此时 $X$ 的每一个开覆盖都是有限的。
\end{example}

\begin{example}\label{example:ch03-half-open-not-compact}
区间 $(0,1]$ 不是紧致的。开覆盖
\begin{align*}
\mathcal{A}=\{\,(1/n,1]\mid n\in\mathbb{Z}_{+}\,\}
\end{align*}
就不包含覆盖 $(0,1]$ 的有限子族。同理，区间 $(0,1)$ 也不是紧致的。另一方面，区间 $[0,1]$ 是紧致的。或许读者已经在数学分析中熟悉了这一结论；尽管如此，稍后我们还是要对此加以证明。
\end{example}

一般说来，判定一个空间是否是紧致空间并非总是轻而易举的。首先，我们要证明几个关于如何从已知紧致空间出发去构造新的紧致空间的一般性定理。下一节，我们再来证明一些特殊空间的紧致性。这些空间包括实直线上所有的闭区间和 $\mathbb{R}^{n}$ 中所有的有界闭子集。

我们先来证明关于子空间的一些结论。设 $Y$ 是 $X$ 的一个子空间，$\mathcal{A}$ 是 $X$ 的一个子集族，如果它的成员之并包含着 $Y$，则称 $\mathcal{A}$ \textbf{覆盖}（cover）$Y$。

\begin{lemma}\label{lemma:ch03-compact-subspace-cover}
设 $Y$ 是 $X$ 的一个子空间。那么，$Y$ 是紧致的当且仅当由 $X$ 的开集所组成的 $Y$ 的每一个覆盖都包含着一个覆盖 $Y$ 的有限子族。
\end{lemma}

\begin{proof}
假设 $Y$ 是紧致的，并且 $\mathcal{A}=\{A_{\alpha}\}_{\alpha\in J}$ 是由 $X$ 的开集所组成的 $Y$ 的一个覆盖。那么族
\begin{align*}
\{\,A_{\alpha}\cap Y\mid\alpha\in J\,\}
\end{align*}
是由 $Y$ 中开集所组成的 $Y$ 的一个覆盖，因此有一个有限子族
\begin{align*}
\{A_{\alpha_{1}}\cap Y,\cdots,A_{\alpha_{m}}\cap Y\}
\end{align*}
覆盖 $Y$。于是 $\{A_{\alpha_{1}},\cdots,A_{\alpha_{m}}\}$ 就是 $\mathcal{A}$ 的一个覆盖 $Y$ 的子族。

反之，假定假设条件成立，我们证明 $Y$ 是紧致的。设 $\mathcal{A}'=\{A_{\alpha}'\}$ 是由 $Y$ 中的开集所构成的 $Y$ 的一个覆盖，对于每一个 $\alpha$，选取 $X$ 中的开集 $A_{\alpha}$ 使得
\begin{align*}
A_{\alpha}'=A_{\alpha}\cap Y.
\end{align*}
从而，族 $\mathcal{A}=\{A_{\alpha}\}$ 是由 $X$ 中开集构成的 $Y$ 的一个覆盖。根据假设，它有一个有限子族 $\{A_{\alpha_{1}},\cdots,A_{\alpha_{m}}\}$ 覆盖 $Y$。于是 $\{A_{\alpha_{1}}',\cdots,A_{\alpha_{m}}'\}$ 便是 $\mathcal{A}'$ 的一个覆盖 $Y$ 的子族。
\end{proof}

\begin{theorem}\label{theorem:ch03-closed-subspace-compact}
紧致空间的每一个闭子集都是紧致的。
\end{theorem}

\begin{proof}
设 $Y$ 是紧致空间 $X$ 的一个闭子集。任意给定由 $X$ 的开集组成的 $Y$ 的一个覆盖 $\mathcal{A}$，将 $\mathcal{A}$ 添加一个开集 $X\setminus Y$ 便构成 $X$ 的一个开覆盖 $\mathcal{B}$，即
\begin{align*}
\mathcal{B}=\mathcal{A}\cup\{X\setminus Y\}.
\end{align*}
$\mathcal{B}$ 的一个有限子族便覆盖了 $X$。如果这个有限子族含有集合 $X\setminus Y$，那么就去掉 $X\setminus Y$。如果没有 $X\setminus Y$，则不再变动。由此得到的族就是 $\mathcal{A}$ 的覆盖 $Y$ 的一个有限子族。
\end{proof}

\begin{theorem}\label{theorem:ch03-compact-subspace-hausdorff-closed}
Hausdorff 空间的每一个紧致子空间都是闭的。
\end{theorem}

\begin{proof}
设 $Y$ 是一个 Hausdorff 空间 $X$ 的一个紧致子空间。我们来证明 $X\setminus Y$ 是开的，从而 $Y$ 是闭的。

令 $x_{0}$ 为 $X\setminus Y$ 的一个点。我们证明存在一个 $x_{0}$ 的与 $Y$ 无交的邻域。对于 $Y$ 中每一个点 $y$，分别选取 $x_{0}$ 和 $y$ 的无交的邻域 $U_{y}$ 和 $V_{y}$（应用 Hausdorff 条件）。于是，$X$ 的开集族 $\{V_{y}\mid y\in Y\}$ 就是 $Y$ 的一个覆盖。因此，就有族中有限多个成员 $V_{y_{1}},\cdots,V_{y_{n}}$ 覆盖 $Y$。开集
\begin{align*}
V=V_{y_{1}}\cup\cdots\cup V_{y_{n}}
\end{align*}
包含 $Y$，并且它与由 $x_{0}$ 的一些相应的邻域取交而构成的开集
\begin{align*}
U=U_{y_{1}}\cap\cdots\cap U_{y_{n}}
\end{align*}
无交。因为，如果 $z$ 是 $V$ 中一点，则对于某一个 $i$ 有 $z\in V_{y_{i}}$，因此 $z\notin U_{y_{i}}$，从而 $z\notin U$。参见\cref{figure:ch03-separate-point-compact}。那么，$U$ 便是 $x_{0}$ 的一个邻域，与 $Y$ 无交。证明完成。
\end{proof}

\par
\begin{figure}[htbp]
\centering
\includegraphics[width=0.55\textwidth]{ch03-fig26-1.png}
\caption{在 Hausdorff 空间中分离点与紧致子空间}\label{figure:ch03-separate-point-compact}
\end{figure}

在前面这个定理的证明中论证的结论将来还会用到，我们在这里重新陈述一次以备引用：

\begin{lemma}\label{lemma:ch03-separate-point-and-compact}
设 $Y$ 是一个 Hausdorff 空间 $X$ 的一个紧致子空间，$x_{0}$ 不属于 $Y$，则存在 $X$ 中的两个无交的开集 $U$ 和 $V$，它们分别包含 $x_{0}$ 和 $Y$。
\end{lemma}

\begin{example}\label{example:ch03-intervals-compact}
一旦我们证明了 $\mathbb{R}$ 中的区间 $[a,b]$ 是紧致的，根据\cref{theorem:ch03-closed-subspace-compact}就可推出 $[a,b]$ 中任何闭子空间都是紧致的。另一方面，从\cref{theorem:ch03-compact-subspace-hausdorff-closed}推出 $\mathbb{R}$ 中区间 $(a,b]$ 和 $(a,b)$ 不可能是紧致的（这是我们已经知道的），因为它们不是 Hausdorff 空间 $\mathbb{R}$ 中的闭集。
\end{example}

\begin{example}\label{example:ch03-finite-complement-compact}
在\cref{theorem:ch03-compact-subspace-hausdorff-closed}的假设中 Hausdorff 条件是必要的，比如，考虑实直线上的有限补拓扑。就这个拓扑而言，$\mathbb{R}$ 的真子集中只有有限集才是闭的。读者可以验证：关于有限补拓扑，$\mathbb{R}$ 中每一个子集都是紧致的。
\end{example}

\begin{theorem}\label{theorem:ch03-continuous-image-compact}
紧致空间的连续像是紧致的。
\end{theorem}

\begin{proof}
设 $f\colon X\to Y$ 连续，$X$ 是紧致的。$\mathcal{A}$ 是由 $Y$ 中开集所构成的 $f(X)$ 的一个覆盖。则集族
\begin{align*}
\{\,f^{-1}(A)\mid A\in\mathcal{A}\,\}
\end{align*}
是 $X$ 的一个覆盖。因为 $f$ 是连续的，这些集合都是 $X$ 中开集。因此它们中有有限多个，比如说
\begin{align*}
f^{-1}(A_{1}),\cdots,f^{-1}(A_{n}),
\end{align*}
便覆盖 $X$。于是 $A_{1},\cdots,A_{n}$ 覆盖 $f(X)$。
\end{proof}

这个定理的重要应用之一是验证一个映射是否是同胚。

\begin{theorem}\label{theorem:ch03-compact-hausdorff-homeo}
设 $f\colon X\to Y$ 是一个连续的一一映射。若 $X$ 是紧致的，并且 $Y$ 是 Hausdorff 的，则 $f$ 是一个同胚。
\end{theorem}

\begin{proof}
我们证明在映射 $f$ 下，$X$ 中闭集的像是 $Y$ 中的闭集。由此便得到映射 $f^{-1}$ 的连续性。设 $A$ 为 $X$ 中闭集，那么根据\cref{theorem:ch03-closed-subspace-compact}得知 $A$ 是紧致的。因此按照刚才证明的那个定理可见 $f(A)$ 是紧致的。因为 $Y$ 是一个 Hausdorff 空间，根据\cref{theorem:ch03-compact-subspace-hausdorff-closed}可见 $f(A)$ 在 $Y$ 中是闭的。
\end{proof}

\begin{theorem}\label{theorem:ch03-finite-product-compact}
有限多个紧致空间的积是紧致的。
\end{theorem}

\begin{proof}
我们来证明两个紧致空间的积是紧致的。对于任意有限积的情形，只要用归纳法便可得到。

第一步。给定空间 $X$ 和 $Y$，其中 $Y$ 是紧致的。又设 $x_{0}$ 是 $X$ 的一个点，$N$ 是 $X\times Y$ 中包含``薄片''$x_{0}\times Y$ 的一个开集。我们证明以下结论：$X$ 中存在 $x_{0}$ 的一个邻域 $W$，使得 $N$ 包含着集合 $W\times Y$。集合 $W\times Y$ 通常称为关于 $x_{0}\times Y$ 的一个\textbf{管子}（tube）。

首先含在 $N$ 中（关于 $X\times Y$ 的拓扑）的那些基元素 $U\times V$ 覆盖 $x_{0}\times Y$。作为一个同胚于 $Y$ 的空间，$x_{0}\times Y$ 是紧致的。因此我们可用有限多个这样的基元素
\begin{align*}
U_{1}\times V_{1},\cdots,U_{n}\times V_{n}
\end{align*}
覆盖 $x_{0}\times Y$（假定每一个基元素 $U_{i}\times V_{i}$ 都与 $x_{0}\times Y$ 相交，如果不然的话，这个基元素就是多余的。我们从这个集的有限族中去掉它之后仍然是 $x_{0}\times Y$ 的一个覆盖）。令
\begin{align*}
W=U_{1}\cap\cdots\cap U_{n}.
\end{align*}
集合 $W$ 是开的，又因为每一个集合 $U_{i}\times V_{i}$ 都与 $x_{0}\times Y$ 相交，所以 $x_{0}$ 属于 $W$。

我们可以断言被选出来覆盖薄片 $x_{0}\times Y$ 的这些集合 $U_{i}\times V_{i}$ 也覆盖着管子 $W\times Y$。设 $x\times y$ 是 $W\times Y$ 的一个点，考虑在薄片 $x_{0}\times Y$ 上与这个点具有相同纵坐标的点 $x_{0}\times y$。对于某一个 $i$ 有 $x_{0}\times y$ 属于 $U_{i}\times V_{i}$，所以 $y\in V_{i}$。而对于每一个 $j$ 有 $x\in U_{j}$（因为 $x\in W$）。于是有 $x\times y\in U_{i}\times V_{i}$。

由于所有这些集合 $U_{i}\times V_{i}$ 都在 $N$ 中，并且它们覆盖 $W\times Y$，于是管子 $W\times Y$ 也包含在 $N$ 中。参见\cref{figure:ch03-tube}。

第二步。现在我们来完成定理证明。设 $X$ 和 $Y$ 都是紧致空间。$\mathcal{A}$ 是 $X\times Y$ 的一个开覆盖。给定 $x_{0}\in X$，薄片 $x_{0}\times Y$ 是紧致的，因此可以用 $\mathcal{A}$ 中有限多个成员 $A_{1},\cdots,A_{m}$ 覆盖它。它们的并 $N=A_{1}\cup\cdots\cup A_{m}$ 是包含着 $x_{0}\times Y$ 的一个开集。根据第一步，开集 $N$ 包含着包含 $x_{0}\times Y$ 的一个管子 $W\times Y$，其中 $W$ 是 $X$ 的一个开集。因此 $W\times Y$ 也被 $\mathcal{A}$ 中的有限多个成员 $A_{1},\cdots,A_{m}$ 所覆盖。

于是，对于 $X$ 中的每一个 $x$，我们可以选择 $x$ 的邻域 $W_{x}$，使得管子 $W_{x}\times Y$ 能被 $\mathcal{A}$ 中有限多个成员所覆盖。所有这些邻域 $W_{x}$ 组成 $X$ 的一个开覆盖。因此，根据 $X$ 的紧致性，存在有限子族
\begin{align*}
\{W_{1},\cdots,W_{k}\}
\end{align*}
覆盖 $X$。这些管子
\begin{align*}
W_{1}\times Y,\cdots,W_{k}\times Y
\end{align*}
的并便是 $X\times Y$。由于每一个管子可以被 $\mathcal{A}$ 中有限多个成员覆盖，所以 $X\times Y$ 也可以被 $\mathcal{A}$ 中有限多个成员覆盖。
\end{proof}

\par
\begin{figure}[htbp]
\centering
\includegraphics[width=0.55\textwidth]{ch03-fig26-2.png}
\caption{管子引理}\label{figure:ch03-tube}
\end{figure}

上面这个证明中第一步的结论以后还要用到，我们把它重新陈述于此以备引用：

\begin{lemma}[管状引理]\label{lemma:ch03-tube-lemma}
考虑积空间 $X\times Y$，其中 $Y$ 是紧致的。如果 $N$ 是 $X\times Y$ 中包含着薄片 $x_{0}\times Y$ 的一个开集，则 $N$ 必包含着关于 $x_{0}\times Y$ 的某一个管子 $W\times Y$，其中 $W$ 是 $x_{0}$ 在 $X$ 中的一个邻域。
\end{lemma}

\begin{example}\label{example:ch03-tube-fails}
如果 $Y$ 不是紧致的，则管状引理不一定正确。比如，设 $Y$ 是 $\mathbb{R}^{2}$ 中的 $y$ 轴，并且设
\begin{align*}
N=\{\,x\times y\mid |x|<1/(y^{2}+1)\,\}.
\end{align*}
则 $N$ 是包含集合 $0\times\mathbb{R}$ 的一个开集，但它不包含关于 $0\times\mathbb{R}$ 的任何管子。参见\cref{figure:ch03-tube-fails}。
\end{example}

\par
\begin{figure}[htbp]
\centering
\includegraphics[width=0.35\textwidth]{ch03-fig26-3.png}
\caption{管状引理失效的例子}\label{figure:ch03-tube-fails}
\end{figure}

很自然地，与之相关的一个问题就是：无穷多个紧致空间的积是紧致的吗？我们期望有肯定的回答，事实也的确如此。这个结果是如此重要（也是相当困难的），以至于要用给出证明的人的名字来为其命名，称之为 Tychonoff 定理。

在证明连通空间的笛卡儿积仍然连通时，我们首先对有限积的情形给予证明，并且由此导出一般情形。然而，为了证明紧致空间还是紧致的，却无法根据有限积情形下的结论推出无限积情形下的结论。无限的情形则要求用一种全新的方法，证明是相当困难的。鉴于证明的难度较大，同时也是为了避免干扰本章讨论的主线，我们决定将其推后。然而，如果读者想尽快知道的话，也可以在本节之后立刻去看它的证明（第 37 节），这不会对内容的连续性有什么影响。

最后，关于空间的紧致性还有一个判别准则，它是用闭集而不是开集来阐述的。乍看起来，这既不自然又不便于使用，而事实上，在许多场合下它是很有用的。我们首先给出一个定义。

\begin{definition}\label{definition:ch03-fip}
$X$ 的一个子集族 $\mathcal{C}$ 称为具有\textbf{有限交性质}（finite intersection property），如果 $\mathcal{C}$ 的任何一个有限子族
\begin{align*}
\{C_{1},\cdots,C_{n}\}
\end{align*}
的交 $C_{1}\cap\cdots\cap C_{n}$ 是非空的。
\end{definition}

\begin{theorem}\label{theorem:ch03-fip-compact}
设 $X$ 是一个拓扑空间。则 $X$ 是紧致的当且仅当 $X$ 中具有有限交性质的每一个闭集族 $\mathcal{C}$，它的所有成员的交 $\bigcap_{C\in\mathcal{C}}C$ 是非空的。
\end{theorem}

\begin{proof}
给定 $X$ 的一个子集族 $\mathcal{A}$，则
\begin{align*}
\mathcal{C}=\{\,X\setminus A\mid A\in\mathcal{A}\,\}
\end{align*}
是它们的补所构成的族。于是我们有下列论断：
\begin{enumerate}
\item $\mathcal{A}$ 是开集族当且仅当 $\mathcal{C}$ 是闭集族.
\item 集族 $\mathcal{A}$ 覆盖 $X$ 当且仅当 $\mathcal{C}$ 的所有成员的交 $\bigcap_{C\in\mathcal{C}}C$ 是空集.
\item $\mathcal{A}$ 的有限子族 $\{A_{1},\cdots,A_{n}\}$ 覆盖 $X$ 当且仅当 $\mathcal{C}$ 中相应成员 $C_{i}=X\setminus A_{i}$ 的交是空集.
\end{enumerate}
论断 (1) 是显然的，而论断 (2) 和论断 (3) 可由 De Morgan 法则
\begin{align*}
X-\Big(\bigcup_{\alpha\in J}A_{\alpha}\Big)=\bigcap_{\alpha\in J}(X-A_{\alpha})
\end{align*}
推出。

现在用两个简单的步骤进行证明：先考虑定理的逆否命题，然后取集合的补！

$X$ 是紧致的这句话等于说：``对 $X$ 的任何开子集族 $\mathcal{A}$，如果 $\mathcal{A}$ 覆盖 $X$，则有 $\mathcal{A}$ 的某一个有限子族覆盖 $X$。''这个论断又等价于它的逆否命题：``对任意开集族 $\mathcal{A}$，如果 $\mathcal{A}$ 中没有有限的子族覆盖 $X$，则 $\mathcal{A}$ 不能覆盖 $X$。''与前面相仿，令 $\mathcal{C}$ 为集族 $\{X\setminus A\mid A\in\mathcal{A}\}$，再应用论断 (1)--(3)，我们看到这个论断等价于``对于任意闭集族 $\mathcal{C}$，如果 $\mathcal{C}$ 中任何有限个成员的交非空，则 $\mathcal{C}$ 的所有成员的交非空。''这恰好就是我们定理的条件。
\end{proof}

当我们有紧致空间 $X$ 闭集的一个\textbf{套序列}（nested sequence）$C_{1}\supset C_{2}\supset\cdots\supset C_{n}\supset C_{n+1}\cdots$ 时，便得到定理的一种特殊情况。容易证明，如果每一个 $C_{n}$ 是非空的，那么集族 $\mathcal{C}=\{C_{n}\}_{n\in\mathbb{Z}_{+}}$ 自然具有有限交性质。于是交
\begin{align*}
\bigcap_{n\in\mathbb{Z}_{+}}C_{n}
\end{align*}
非空。

我们在下一节中证明实数集的不可数性，在第 5 章中证明 Tychonoff 定理以及在第 8 章中证明 Baire 范畴定理的时候，都要用到这个紧致性的闭集判别准则。

\subsection*{习题}

\begin{enumerate}
\item (a) 设 $\mathscr{T}$ 和 $\mathscr{T}'$ 是 $X$ 的两个拓扑，并且 $\mathscr{T}'\supset\mathscr{T}$。$X$ 关于其中哪一个拓扑是紧致的可以推出对另一个拓扑 $X$ 是紧致的？
\item (b) 证明：若 $X$ 关于 $\mathscr{T}$ 和 $\mathscr{T}'$ 都是紧致的 Hausdorff 空间，则或者 $\mathscr{T}$ 和 $\mathscr{T}'$ 相等，或者它们不能比较。
\item (a) 证明：相对于有限补拓扑而言，$\mathbb{R}$ 的任何子集都是紧致的。
\item (b) 由 $\mathbb{R}$ 的满足 $\mathbb{R}\setminus A$ 是一个可数字集或者是整个 $\mathbb{R}$ 的所有子集 $A$ 构成的 $\mathbb{R}$ 的拓扑，相对于这个拓扑 $[0,1]$ 是紧致子空间吗？
\item 证明：紧致子空间的有限并是紧致的。
\item 证明：度量空间的每一个紧致子空间，对于给定的度量而言是有界的并且是闭的。找一个度量空间，在它里面并不是每一个有界闭子集都是紧致的。
\item 设 $A$ 和 $B$ 是一个 Hausdorff 空间中的两个无交的紧致子空间。证明存在分别包含 $A$ 和 $B$ 的无交的开集 $U$ 和 $V$。
\item 证明：若 $f\colon X\to Y$ 是连续的，其中 $X$ 是紧致的，$Y$ 是 Hausdorff 的，则 $f$ 是闭映射（也就是 $f$ 将闭集映为闭集）。
\item 证明：若 $Y$ 是紧致的，则投射 $\pi_{1}\colon X\times Y\to X$ 是闭映射。
\item \textbf{定理}\quad 设 $f\colon X\to Y$，$Y$ 是紧致的 Hausdorff 空间。则 $f$ 连续当且仅当 $f$ 的\textbf{图形}（graph）
\begin{align*}
G_{f}=\{\,x\times f(x)\mid x\in X\,\}
\end{align*}
是 $X\times Y$ 的闭集。[提示：若 $G_{f}$ 是闭的，$V$ 是 $f(x_{0})$ 的一个邻域，则 $G_{f}$ 与 $X\times(Y\setminus V)$ 的交为闭集。应用习题 7 的结论。]
\item 下面是管状引理的推广：

\textbf{定理}\quad 设 $A$ 和 $B$ 分别是 $X$ 和 $Y$ 的子集，$N$ 是 $X\times Y$ 中包含 $A\times B$ 的一个开集。若 $A$ 和 $B$ 都是紧致的，则在 $X$ 和 $Y$ 中分别存在开集 $U$ 和 $V$，使得
\begin{align*}
A\times B\subset U\times V\subset N.
\end{align*}
\item (a) 证明以下一致极限定理的部分逆：

\textbf{定理}\quad 设 $f_{n}\colon X\to\mathbb{R}$ 是连续函数的一个序列，并且对于每一个 $x\in X$ 有 $f_{n}(x)\to f(x)$。若 $f$ 连续，$f_{n}$ 单调上升并且 $X$ 是紧致的，则这个收敛是一致收敛的。[若对于所有 $n$ 和 $x$ 有 $f_{n}(x)\leqslant f_{n+1}(x)$，则称 $f_{n}$ 是单调上升的。]

(b) 举例说明，如果去掉 $X$ 是紧致的这个条件，或去掉序列是单调的这个条件，则定理不成立。[提示：见第 21 节习题。]
\item \textbf{定理}\quad 设 $X$ 是紧致的 Hausdorff 空间，$\mathcal{A}$ 是 $X$ 的闭的连通子集的一个族，并且 $\mathcal{A}$ 在真包含关系之下是全序的，则
\begin{align*}
Y=\bigcap_{A\in\mathcal{A}}A
\end{align*}
是连通的。[提示：若 $C\cup D$ 是 $Y$ 的一个分割，在 $X$ 中找出无交的开集 $U$ 和 $V$，它们分别包含 $C$ 和 $D$，然后证明
\begin{align*}
\bigcap_{A\in\mathcal{A}}(A\setminus(U\cup V))
\end{align*}
不是空集。]
\item 设 $p\colon X\to Y$ 是一个闭连续满射，对于每一个 $y$，$p^{-1}(\{y\})$ 是一个紧致空间。（这样的映射也称之为\textbf{完备映射}（perfect map）。）证明：若 $Y$ 是紧致的，则 $X$ 是紧致的。[提示：若 $U$ 是包含 $p^{-1}(\{y\})$ 的一个开集，那么存在一个 $y$ 的邻域 $W$，使得 $p^{-1}(W)$ 被 $U$ 所包含。]
\item 设 $G$ 是一个拓扑群。
\begin{enumerate}
\item 设 $A$ 和 $B$ 是 $G$ 的一个子集，如果在 $G$ 中 $A$ 是闭的并且 $B$ 是紧致的，则 $A\cdot B$ 在 $G$ 中是闭的。[提示：若 $c$ 不在 $A\cdot B$ 中，找一个 $c$ 的邻域 $W$，使得 $W\cdot B^{-1}$ 与 $A$ 无交。]
\item 设 $H$ 是 $G$ 的一个子群，$p\colon G\to G/H$ 是一个商映射。若 $H$ 是紧致的，证明 $p$ 是一个闭映射。
\item 设 $H$ 是 $G$ 的一个紧致子群。证明：若 $G/H$ 是紧致的，则 $G$ 是紧致的。
\end{enumerate}
\end{enumerate}